\section{Algorithms whose state is a forest}\label{sec:forests}

This is the section that decides the claim. Four bodies of work sit between
\FSC{} and any statement of the form ``an algorithm whose state is a forest on
the precedence graph''. The first of them is sixty years old and comes from
mining, which is where this project's own instances come from; it was missing
from the earlier drafts of this review, and it is the most uncomfortable of the
four.

\subsection{Normalized trees: Lerchs--Grossmann, 1965}\label{sec:lg}

\citet{LerchsGrossmann1965} solve maximum closure with a graph algorithm, not
with a flow algorithm, and the object it maintains is a forest. In the exposition
of \citet{Hochbaum2001}, from which the following is taken verbatim in substance:
the algorithm maintains a spanning tree rooted at an added vertex $r$ in the
extended graph, called a \emph{normalized tree}; each branch $T_q$ hanging from
the root carries its \emph{mass} $M_q = b(T_q)$, the sum of the vertex weights it
contains; a downward arc is \emph{strong} if the mass it supports is positive;
the tree is normalized when its only strong arcs are the ones adjacent to $r$;
and, dropping the root arcs, Hochbaum writes that ``we also refer to the forest
$(V, A_T)$ in $G$ as a normalized tree''. The invariant is that the union of the
strong branches is a maximum closed set of the forest, and the weight of the
strong vertices bounds the maximum closure value from above.

An iteration is a pivot in everything but the name. Quoting \citet{Hochbaum2001}
again: a \emph{merger} consists of ``adding an arc $(s,w)$ from a strong branch
$S$ to a weak branch $W$ and removing the arc from the root of the branch $S$ to
the root $r$'', after which the branch hangs from the strong vertex $s$ and the
supported masses are updated. One arc enters, one arc leaves, one component
loses its root, and the aggregates on the affected path are recomputed.
\citet{ZhaoKim1992} is a variant that differs in how the blocks are regrouped
after a violation, which is the split side of the same operation, and
\citet{GianniniEtAl1991} is the flow-based implementation that competed with it.

Two further facts from \citet{Hochbaum2001} matter here. First, the normalized
tree carries enough information to construct a feasible flow in linear time, so
the primal--dual correspondence between the forest state and a flow certificate
is known. Second, he gives a parametric version of the Lerchs--Grossmann
algorithm with the same complexity as a single run --- which makes it a
\emph{competitor} on the task \FS{} performs, not merely an antecedent.
\citet{PicardSmith2004} close the loop on the objective as well: they compute
intermediate open-pit contours maximising a benefit--cost ratio by parametric
maximum flow.

\paragraph{What this costs the claim.} A great deal, and more than the recent
work does. The correspondence with \FS{} is not thematic but term by term: a
branch aggregate $b(T_q)$ against the reduced cost $p(T) - \rho\,w(T)$; a strong
arc against a positive reduced cost; the union of the strong branches against
the support of the current basic solution; a merger against an entering arc with
its leaving root arc. Setting $b_i = p_i - \rho w_i$ makes the Lerchs--Grossmann
strong/weak test literally the sign of the \FS{} reduced cost. So the forest
state, the branch aggregates, the sign test on an aggregate and the merge
operation are not new relative to the 2025--2026 work: they are not new relative
to 1965. What Lerchs--Grossmann do \emph{not} do is run a simplex --- there is no
basis in the LP sense, no ratio test, no dual multiplier as an LP object --- and
they solve the closure problem at fixed weights rather than the ratio LP.
That distinction is now the whole of the surviving claim, and
\cref{sec:open}~Q5 says what has to be proved to keep it.

\subsection{Network simplex with a side constraint}

That a basis of a network LP is a spanning tree, that adding a side row makes it
a spanning structure plus one extra element, and that pivots are described as
operations on that structure, is textbook material
\citep{AhujaEtAl1993,Schrijver2003}. The paradigm ``simplex on an incidence
matrix with one side row'' is therefore not available as a contribution, and any
formulation that suggests otherwise must be rewritten.

The specific case of one side row has its own line. \citet{ChenSaigal1977} give
a revised simplex variant for capacitated network flow with additional linear
constraints in which the working basis has the size of the side system rather
than of the network. \citet{Caliskan2011} specialises the network simplex to
the constrained maximum flow problem. \citet{HolzhauserEtAl2017} give a fully
combinatorial network simplex for the budget-constrained minimum cost flow
problem, with two kinds of vertex potentials, three kinds of reduced costs,
optimality criteria proved directly on the network and an explicit anti-cycling
treatment. These are the right block to read, but they are one
transposition away from \eqref{eq:class}: their variables are arcs and their
rows are vertices, whereas here the variables are vertices and the rows are arcs
(\cref{sec:class}(c)). The two formulations are linear programming duals, so the
machinery transports and the differences have to be argued rather than assumed
(\cref{sec:open}, Q3) --- but the resemblance is a duality, not an identity, and
an earlier draft overstated it.

\subsection{Yang: a specialised simplex on rooted spanning forests}

\citet{Yang2026} gives, for budget-constrained total-variation-regularised
linear programs on graphs, a characterisation of basic solutions in terms of
rooted spanning forests with orientation, an interpretation of the simplex
method as operations on those forests, and an accelerated implementation
compared against a commercial solver. The abstract states the first two points
directly. Total variation on a graph is a difference system on vertex variables
with a budget row, so his program and \eqref{eq:class} are in the same
structural family, and not merely in the same spirit: this is the closest match
in formulation type among all the work reviewed here.

The overlap with \FS{} is substantial and was underestimated in an earlier
draft. Yang's setting has a budget row and a graph-structured constraint matrix;
when the budget slack is nonbasic he obtains a component without a root, which
plays the role of the rootless positive component here; he writes the basic
variables explicitly in terms of the nonbasic ones, derives reduced costs for
root and edge variables, states optimality in terms of tree and subtree shifts,
reads pivots as merges, splits and shifts, uses blocking edges, and updates
subtree aggregates locally.

Consequently the following claim, made in an earlier draft, is withdrawn: that
no work carries revised-simplex algebra to this level of explicitness on a
forest structure. Yang does exactly that, on a different linear program.

\subsection{Wuille and the spanning-forest linearization line}

\citet{Wuille2025} sets up, for Bitcoin transaction cluster linearization, the
program \(\max \sum_i f_i t_i\) subject to \(\sum_i s_i t_i = 1\), the dependency
inequalities and \(t \ge 0\): the same normalised program \FS{} solves, up to
orientation convention and application. The post states that the method
maintains \(n-1\) free variables, that the active dependencies form a spanning
forest, that within every chunk there is at most one free vertex variable, that
all chunks except one contain such a variable and are therefore excluded, and it
describes the pivot operations as merging and splitting chunks. These were
verified against the post itself.

\citet{MollakhaniEtAl2026} formalise that reading: the paper presents the
equivalent linear programming formulation and develops the Spanning Forest
Linearization algorithm, motivated, in the authors' words, by structural
properties of basic feasible solutions of the simplex method represented through
spanning forests. It then leaves the exact simplex: it stops maintaining the
identity of the root and a single distinguished included component, compares
chunks locally, and optimises the linearization directly.

This is the closest antecedent in absolute terms, and it removes several items
from the candidate contribution list: the normalised program itself, the count
of \(n-1\) nonbasic variables, the forest of tight slacks, the ``at most one
free vertex variable per component'', the single positive component, the
connectivity of the support, and the four qualitative pivot effects. None of
these may be presented as new.

\subsection{What the four leave open}

Lerchs--Grossmann maintain a forest on the closure graph and test the sign of
branch aggregates, but do not run a simplex and do not solve the ratio LP.
\citet{Yang2026} develops the exact simplex, but for a different program. Wuille
and Mollakhani et al.\ work on the same program, but leave the exact simplex in
favour of a direct linearization algorithm. \citet{HolzhauserEtAl2017} develops a
combinatorial simplex for a network with one side row, but on minimum cost flow.
The intersection of all four --- the ordinary primal simplex, developed exactly,
on the maximum-ratio closure program --- is where whatever remains of the
contribution has to live, and \cref{sec:positioning} states how little that is
and what would have to be shown to make it stand.
